\chapter{量子角动量}\label{ch:b3:quantum-angular-momentum}

将射电望远镜对准一片暗黑的天空，并将其调谐至
\qty{115}{GHz}：一条明亮而尖锐的谱线随即出现——这是
一氧化碳分子，温度比绝对零度高十开尔文，正从
其\emph{转动}态阶梯上跃下一级。它们的翻滚运动，与
量子力学中一切转动的情形一样，受到双重量子化：
角动量的大小只能取
$\sqrt{l(l+1)}\,\hbar$，而其沿任一选定轴的分量只能取
$m\hbar$。本章将推导这一双重量子化——一是代数地，
即由\omterm{def:b3:quantum-formalism:commutator}{对易子}出发，这些\omterm{def:b3:quantum-formalism:commutator}{对易子}继承自
\cref{ch:b3:hamiltonian-mechanics}中的\omterm{def:b3:hamiltonian-mechanics:poisson}{泊松括号}；另一是具体地，
即作为塑造每个原子轨道的\omterm{prop:b3:quantum-angular-momentum:harmonics}{球谐函数}。代数
所给出的结果比我们所求的更多：它允许取半整数值，
而任何轨道运动都无法实现它们，自然界将在两章之后用自旋填补这一空位。
途中我们还将遇到分子的
转动阶梯——毫米波谱线，天文学家借以
测量星系中冷气体的质量——以及
角动量最早借以显示自身分裂的磁矩。

\section{转动的代数}

\begin{definition}[角动量算符]\label{def:b3:quantum-angular-momentum:operators}
\index{角动量!量子}
轨道角动量为算符 $\hat{\vect L} =
\hat{\vect r}\wedge\hat{\vect p}$，其分量为 $\hat L_z = \hat
x\hat p_y - \hat y\hat p_x$，其余分量依轮换得到。由 $[\hat x, \hat p_x] =
\iu\hbar$ 出发，有：
\[
[\hat L_x, \hat L_y] = \iu\hbar\,\hat L_z
\quad\text{（及轮换）} , \qquad
[\hat L^2, \hat L_z] = 0 :
\]
各分量彼此不相容——没有任何态能使其中两个
同时取确定值——但总平方与其中任一分量相容：
数对 $(\hat L^2, \hat L_z)$ 是标准的
标记选择。这些\omterm{def:b3:quantum-formalism:commutator}{对易子}恰为 $\iu\hbar$ 乘以
\cref{ex:b3:hamiltonian-mechanics:brackets}中算出的\omterm{def:b3:hamiltonian-mechanics:poisson}{泊松括号}：
这正是狄拉克对应法则的应用。
\end{definition}

\begin{theorem}[由代数本身得出的谱]\label{thm:b3:quantum-angular-momentum:spectrum}
\index{磁量子数}
设 $\hat J_x, \hat J_y, \hat J_z$ 为\emph{任意}三个\omterm{def:b3:quantum-formalism:observable}{厄米
算符}，满足上述对易关系。则
$(\hat J^2, \hat J_z)$ 的联合\omterm{def:b3:quantum-formalism:observable}{本征值}为
\[
\hat J^2:\ j(j+1)\hbar^2 , \qquad
\hat J_z:\ m\hbar , \quad m = -j, -j+1, \dots, +j ,
\]
其中 $j$ 为非负的\emph{整数或半整数}：$2j + 1$ 个
$m$ 值，对每个 $j$。\omterm{def:b3:harmonic-oscillator:ladder}{升降算符} $\hat J_\pm = \hat
J_x \pm \iu\hat J_y$ 使 $m$ 每次改变 $\pm1$，而 $j$ 固定：
\[
\hat J_\pm\ket{j,m} = \hbar\sqrt{j(j+1) - m(m\pm1)}\,\ket{j,m\pm1} .
\]
\end{theorem}

\begin{proof}[部分证明]
$[\hat J_z, \hat J_\pm] = \pm\hbar\hat J_\pm$：作用 $\hat
J_\pm$ 使 $\hat J_z$ 的\omterm{def:b3:quantum-formalism:observable}{本征值}改变 $\pm\hbar$，
在 $\hat J^2$ 保持不变的情况下（它与由 $\hat
J_i$ 构造的一切量对易）。范数 $\|\hat J_\pm\ket{j,m}\|^2 = \hbar^2[j(j+1) -
m(m\pm1)]$（由 $\hat J_\mp\hat J_\pm = \hat J^2 - \hat
J_z^2 \mp \hbar\hat J_z$）必须非负：阶梯须在上方某一
$m_{\max}$ 处终止，满足 $m_{\max}(m_{\max}+1) =
j(j+1)$，即 $m_{\max} = j$，下方则终止于 $m_{\min} = -j$。
从 $-j$ 以单位步长攀升至 $+j$，迫使 $2j$
为非负整数。\omterm{def:b3:quantum-formalism:observable}{本征值}记号 $j(j+1)\hbar^2$ 的写法
由此事后得到证实。
\end{proof}

\begin{remark}[目录中的空位]\label{rem:b3:quantum-angular-momentum:halfinteger}
代数允许 $j = \tfrac12, \tfrac32, \dots$——即
含偶数级台阶的阶梯。我们接下来将说明，轨道运动只用到
整数。然而，自然界并不浪费：电子本身携带 $j = \tfrac12$，
却没有轨道与之对应
（见\cref{ch:b3:spin-two-level}）；此处推导的代数保持不变，
将支撑原子磁性、核自旋以及量子比特。
\end{remark}

\section{轨道角动量：球谐函数}

\begin{proposition}[球谐函数]\label{prop:b3:quantum-angular-momentum:harmonics}
\index{球谐函数}\index{轨道量子数}
在球坐标下 $\hat L_z = -\iu\hbar\,\partial_\varphi$，
且 $(\hat L^2, \hat L_z)$ 在球面上的联合本征函数即
\emph{\omterm{prop:b3:quantum-angular-momentum:harmonics}{球谐函数}} $Y_l^m(\theta, \varphi)$：
\[
\hat L^2\,Y_l^m = l(l+1)\hbar^2\,Y_l^m , \qquad
\hat L_z\,Y_l^m = m\hbar\,Y_l^m ,
\]
其中 $l = 0, 1, 2, \dots$ 为\emph{整数}——由
$\eu^{\iu m\varphi}$ 的单值性要求 $m$ 为整数，进而 $l$ 为整数。
前几个函数（相差归一化常数）为：$Y_0^0 = \text{常数}$（$s$
形，球面）；$Y_1^0 \propto \cos\theta$ 与 $Y_1^{\pm1}
\propto \sin\theta\,\eu^{\pm\iu\varphi}$（$p$ 形，双瓣）；
$Y_2^0 \propto 3\cos^2\theta - 1$ 及其同族函数（$d$ 族）。
宇称：$Y_l^m(-\vect r$ 方向$) = (-1)^l\,Y_l^m$。
\end{proposition}
\admitted

\begin{omfigure}
\begin{tikzpicture}[font=\small, scale=0.94]
  % s orbital
  \begin{scope}
    \draw[->] (0,-1.7) -- (0,1.8);
    \node[anchor=south] at (0,1.8) {$z$};
    \draw[omDef, thick, fill=omDef!18] (0,0) circle (1.1);
    \node at (0,-2.3) {$l = 0$ ($s$): 球面};
  \end{scope}
  % p_z
  \begin{scope}[xshift=4.4cm]
    \draw[->] (0,-1.7) -- (0,1.8);
    \node[anchor=south] at (0,1.8) {$z$};
    \draw[omThm, thick, fill=omThm!18, domain=0:360, smooth, samples=100]
      plot ({1.45*abs(cos(\x))*sin(\x)}, {1.45*abs(cos(\x))*cos(\x)});
    \node at (0,-2.3) {$l = 1$, $m = 0$: $|\cos\theta|$};
  \end{scope}
  % p_+-1
  \begin{scope}[xshift=8.8cm]
    \draw[->] (0,-1.7) -- (0,1.8);
    \node[anchor=south] at (0,1.8) {$z$};
    \draw[omExo, thick, fill=omExo!18, domain=0:360, smooth, samples=100]
      plot ({1.45*abs(sin(\x))*sin(\x)}, {1.45*abs(sin(\x))*cos(\x)});
    \node at (0,-2.3) {$l = 1$, $m = \pm1$: $|\sin\theta|$};
  \end{scope}
  % d_z2
  \begin{scope}[xshift=13.2cm]
    \draw[->] (0,-1.7) -- (0,1.8);
    \node[anchor=south] at (0,1.8) {$z$};
    \draw[omProp, thick, fill=omProp!18, domain=0:360, smooth, samples=140]
      plot ({0.62*abs(3*cos(\x)*cos(\x)-1)*sin(\x)}, {0.62*abs(3*cos(\x)*cos(\x)-1)*cos(\x)});
    \node at (0,-2.3) {$l = 2$, $m = 0$};
  \end{scope}
\end{tikzpicture}

{\small $|Y_l^m(\theta)|$ 的极坐标图（取平面内的截面，该平面
含 $z$ 轴；完整图形为绕该轴旋转所得的图形）。
这些与势能无关的角向骨架，
在\cref{ch:b3:hydrogen-atom}中，它们将塑造每一个原子的角向结构。}
\end{omfigure}

\begin{example}[刚性转子及其能级]\label{ex:b3:quantum-angular-momentum:rotor}
转动惯量为 $I$ 的双原子分子作转动，\omterm{def:b3:hamiltonian-mechanics:hamiltonian}{哈密顿量} $\hat H
= \hat L^2/2I$，能量为
\[
E_J = \frac{\hbar^2}{2I}\,J(J+1) = B\,J(J+1) ,
\qquad J = 0, 1, 2, \dots
\]
每一能级为 $(2J+1)$ 重简并。光子吸收满足 $\Delta J =
\pm1$，因此吸收频率为 $\nu_{J+1\leftarrow J} =
2B(J+1)/h$：一组\emph{等间距}谱线——是转动光谱的特征，
据此一眼即可辨认转动光谱。对
一氧化碳，$B/h = \qty{57.6}{GHz}$：基频谱线位于
\qty{115}{GHz}（$\lambda = \qty{2.6}{mm}$），是毫米射电天文学中的
主力谱线
（\cref{pb:b3:quantum-angular-momentum:1}）。
\end{example}

\section{有心势，补充完整}

\begin{theorem}[任意有心势中的分离变量]\label{thm:b3:quantum-angular-momentum:radial}
\index{径向方程}\index{离心势垒}
对 $V(r)$，定态可取为
\[
\psi(\vect r) = \frac{u(r)}{r}\,Y_l^m(\theta, \varphi) ,
\]
其中 $u$ 满足一维\omterm{prop:b3:schrodinger-three-dimensions:swave}{径向方程}
\[
-\frac{\hbar^2}{2m}\,u'' + \Big[V(r) +
\frac{\hbar^2\,l(l+1)}{2mr^2}\Big]u = E\,u , \qquad u(0) = 0 .
\]
角向问题由 $Y_l^m$ 一劳永逸地解决；每个
$l$ 附加一项排斥性的\emph{\omterm{thm:b3:quantum-angular-momentum:radial}{离心势垒}}
$\hbar^2l(l+1)/2mr^2$——是
\cref{ex:b3:lagrangian-mechanics:central}中有效势能的量子版本——且
给定 $l$ 的每一能级均为 $(2l+1)$ 重简并，因为
没有任何有心力能关心 $z$ 轴的取向。$s$ 波
技巧（见\cref{prop:b3:schrodinger-three-dimensions:swave}）即为
本定理中 $l = 0$ 的那一行。
\end{theorem}

\begin{proof}[部分证明]
拉普拉斯算子可分解为 $\Delta = \tfrac1r\partial_r^2(r\,\cdot) -
\hat L^2/\hbar^2r^2$（在本阶段不加证明地承认，它表明 $\hat L^2$
是 $-\hbar^2\Delta$ 的角向部分）。代入 $\psi = (u/r)Y_l^m$，
并利用 $\hat L^2$ 的\omterm{def:b3:quantum-formalism:observable}{本征值}，得到所述方程。
简并性由以下事实推出：$\hat H$ 与全部三个 $\hat
L_i$ 对易，其\omterm{def:b3:harmonic-oscillator:ladder}{升降算符}在不改变能量的前提下改变 $m$。
\end{proof}

\begin{omfigure}
\begin{tikzpicture}
  \begin{axis}[omaxis, axis x line=bottom, axis y line=left,
    width=10.6cm, height=6.0cm, xmin=0, xmax=8, ymin=-1.35, ymax=1.3,
    xlabel={$r$（单位：$a_0$）}, ylabel={$U_{\text{有效}}$（单位：$E_{\text{I}}$）},
    xlabel style={at={(axis description cs:0.5,-0.1)}, anchor=north},
    xtick={2,4,6}, ytick=\empty,
    legend style={font=\scriptsize, at={(0.97,0.95)}, anchor=north east,
      draw=none, fill=none}]
    \addplot[omDef, thick, domain=1.35:8, samples=120] {-2/x};
    \addlegendentry{$l = 0$：纯库仑势}
    \addplot[omThm, thick, domain=0.42:8, samples=160] {-2/x + 1/(x*x)};
    \addlegendentry{$l = 1$}
    \addplot[omExo, thick, domain=0.85:8, samples=160] {-2/x + 3/(x*x)};
    \addlegendentry{$l = 2$}
    \draw[gray, dashed] (axis cs:0,0) -- (axis cs:8,0);
  \end{axis}
\end{tikzpicture}

{\small 库仑问题的有效径向势能：
\omterm{thm:b3:quantum-angular-momentum:radial}{离心势垒} $\hbar^2l(l+1)/2mr^2$ 在
$l \ge 1$ 时把原点隔开。只有 $s$ 态触及原子核——
其影响遍及原子光谱乃至放射性电子俘获。}
\end{omfigure}

\section{磁矩：让角动量可见}

\begin{proposition}[轨道磁矩与塞曼效应]\label{prop:b3:quantum-angular-momentum:magnetic}
\index{玻尔磁子}\index{塞曼效应}
电荷为 $q$、质量为 $m$ 的粒子，具有轨道角动量
$\hat{\vect L}$，则其磁矩为
\[
\hat{\vect\mu} = \frac{q}{2m}\,\hat{\vect L} ;
\]
对电子而言，自然单位是\emph{\omterm{prop:b3:quantum-angular-momentum:magnetic}{玻尔磁子}}
$\mu_{\text{B}} = e\hbar/2m_{\text{e}} = \qty{5.79e-5}{eV/T}$
（参见\cref{exo:b3:schrodinger-three-dimensions:5}）。在
磁场 $B\vect e_z$ 中，能量 $-\hat{\vect\mu}\cdot\vect B$ 使每一能级
移动 $+m\,\mu_{\text{B}}B$（电子电荷为负）：一个
给定 $l$ 的能级分裂为其 $2l + 1$ 个分量——
\emph{\omterm{prop:b3:quantum-angular-momentum:magnetic}{塞曼效应}}，是
$m$ 量子化的首次直接展示，也解释了 $m$ 为何称为\omterm{thm:b3:quantum-angular-momentum:spectrum}{磁量子数}。
在 $B = \qty{1}{T}$ 时分裂为 $\qty{58}{\micro eV}$，
与光学能量相比很小，可以很容易地作为谱线
位移被分辨出来——反过来看，即为磁强计：太阳光谱线的塞曼分裂
可用来测绘太阳黑子的磁场。
\end{proposition}

\begin{proof}[部分证明]
经典地，轨道上的电荷相当于一个电流环：$\mu = IA =
(qv/2\pi r)(\pi r^2) = qL/2m$；算符表述沿用此结果
（它也可由 $\hat{\vect A}$ 的耦合导出，参见
\cref{prop:b3:lagrangian-mechanics:charge}）。能量移动属于
一级微扰理论，此处先行引用，
并在\cref{ch:b3:perturbation-theory}中予以论证。
\end{proof}

\begin{omfigure}
\begin{tikzpicture}[font=\small, >=Stealth]
  \draw[->] (0,-3.0) -- (0,3.15);
  \node[anchor=south] at (0,3.15) {$z$};
  % sphere radius sqrt6
  \pgfmathsetmacro\R{2.449}
  \draw[gray, dashed] (0,0) circle (\R);
  \foreach \m/\c in {2/omDef, 1/omThm, 0/omExo, -1/omThm, -2/omDef} {
    \pgfmathsetmacro\zz{\m}
    \pgfmathsetmacro\rr{sqrt(\R*\R-\m*\m)}
    \draw[\c, thick, ->] (0,0) -- (\rr,\zz);
    \draw[gray!60, thin, dashed] (-0.05,\zz) -- (\rr,\zz);
    \node[anchor=west, \c] at ({\rr+0.1},\zz) {$m = \m$};
  }
  \node[anchor=east, font=\footnotesize] at (-0.15,2.0) {$L_z = m\hbar$};
  \node[anchor=north, font=\footnotesize] at (1.3,-2.75) {$\|\vect L\| = \sqrt{6}\,\hbar$};
\end{tikzpicture}

{\small “空间量子化”（$l = 2$）：角动量
矢量的长度为 $\sqrt{l(l+1)}\,\hbar = \sqrt6\,\hbar$，但
$z$ 轴却只能取五个投影 $m\hbar$——永远达不到其
完整长度：由于各分量彼此不相容，$\vect L$ 永远
不可能恰好沿任一坐标轴。}
\end{omfigure}

\begin{method}[角动量的实际应用]\label{met:b3:quantum-angular-momentum:method}
(1) 用 $(l, m)$ 标记态——抽象代数中则用 $(j, m)$——
且永远不要同时询问两个分量。(2) 矩阵元：
使用升降公式；$\hat L_x = (\hat L_+ + \hat
L_-)/2$。(3) 中心势：引用 $Y_l^m$，仅
求解带\omterm{thm:b3:quantum-angular-momentum:radial}{离心势垒}的\omterm{prop:b3:schrodinger-three-dimensions:swave}{径向方程}。(4) 转动能：
$BJ(J+1)$，谱线位于 $2B(J+1)$——由等间距提取 $B$，进而得到键
长。(5) 磁学问题：角动量换算为磁矩时，以
$\mu_{\text{B}}$ 每 $\hbar$ 计，能量以
$\mu_{\text{B}}B$ 每单位 $m$ 计。
\end{method}

\section{练习}

\begin{exercise}[$\star$]\label{exo:b3:quantum-angular-momentum:1}
(a) 由正则\omterm{def:b3:quantum-formalism:commutator}{对易子}推出 $[\hat L_x, \hat L_y] = \iu\hbar\hat L_z$。
(b) 证明 $[\hat L^2, \hat L_z] = 0$。(c) 证明
$[\hat L_z, \hat x] = \iu\hbar\hat y$：$\hat L_z$ 会
生成什么变换？(d) 三个分量为何没有共同的本征基
——唯有某个特定态例外（指出是哪一个）？
\end{exercise}

\begin{exercise}[$\star$]\label{exo:b3:quantum-angular-momentum:2}
对 $l = 1$，在基 $\{\ket{1,1}, \ket{1,0}, \ket{1,-1}\}$ 中：
(a) 写出 $\hat L_z$ 的矩阵；(b) 用升降公式
写出 $\hat L_+$ 和 $\hat L_-$；(c) 组装 $\hat L_x$，并检验
其\omterm{def:b3:quantum-formalism:observable}{本征值}为 $\hbar, 0, -\hbar$；(d) 对 $L_x =
+\hbar$ 的态沿 $z$ 方向测量，求三种
结果的概率。
\end{exercise}

\begin{exercise}[$\star$]\label{exo:b3:quantum-angular-momentum:3}
一氧化碳：键长 \qty{0.113}{nm}，约化质量
$\qty{1.14e-26}{kg}$。(a) 计算 $I$ 及 $B = \hbar^2/2I$（以 meV 为单位）。
(b) 求 $J = 1 \leftarrow 0$ 谱线的频率和波长。
(c) 求接下来的两条谱线。(d) 天文学家将梳状谱线间距
测量到五位有效数字：她得到的是何种分子物理量，
精度可达多少？
\end{exercise}

\begin{exercise}[$\star$]\label{exo:b3:quantum-angular-momentum:4}
粒子处于 $l = 2$ 的态。(a) 列出测量
$L_z$ 所有可能的结果。(b) 求 $\vect L$ 与
$z$ 轴之间的最小夹角，利用 $\cos\theta = m/\sqrt{l(l+1)}$ 求值。(c) 为什么
该夹角永远不可能为零？(d) 证明当 $l \to \infty$ 时，
最小夹角趋于零：经典矢量得以恢复。
\end{exercise}

\begin{exercise}[$\star\star$]\label{exo:b3:quantum-angular-momentum:5}
$-\hbar^2\Delta$ 的角向部分为 $\hat L^2$，其中
\[
\hat L^2 = -\hbar^2\Big[\frac{1}{\sin\theta}\,\partial_\theta
(\sin\theta\,\partial_\theta) +
\frac{1}{\sin^2\theta}\,\partial_\varphi^2\Big] .
\]
(a) 验证 $Y_1^0 \propto \cos\theta$ 的 $\hat L^2$ \omterm{def:b3:quantum-formalism:observable}{本征值}
为 $2\hbar^2$。(b) 验证 $Y_1^{\pm1} \propto \sin\theta\,
\eu^{\pm\iu\varphi}$ 同样成立，并给出其 $\hat L_z$ \omterm{def:b3:quantum-formalism:observable}{本征值}。(c)
检验其宇称。(d) 为何必有 $\braket{Y_1^0}{Y_1^1} = 0$，
而无需计算任何积分？
\end{exercise}

\begin{exercise}[$\star\star$]\label{exo:b3:quantum-angular-momentum:6}
哪条转动谱线最亮？能级 $J$ 的布居数
$\propto (2J+1)\,\eu^{-BJ(J+1)/k_{\text{B}}T}$。(a) 解释这两个
因子。(b) 证明极大值位于 $J_{\max} \approx
\sqrt{k_{\text{B}}T/2B} - \tfrac12$ 附近。(c) 对
温度为 \qty{10}{K}（暗云）和 \qty{300}{K} 的 CO，哪些谱线占主导？(d)
反推：观测到的 CO 转动谱线序列在 $J = 7$ 处达峰，
说明温度为多少？
\end{exercise}

\begin{exercise}[$\star\star$]\label{exo:b3:quantum-angular-momentum:7}
正常\omterm{prop:b3:quantum-angular-momentum:magnetic}{塞曼效应}。一条波长为 \qty{500}{nm} 的谱线来自
$l = 2 \to l = 1$ 的跃迁，置于磁场 $B = \qty{2}{T}$ 中；忽略
自旋（适用于特殊的“单态”）。(a) 画出分裂后的
能级草图。(b) 利用选择定则 $\Delta m = 0, \pm1$，证明仅
出现\emph{三条}谱线，位置分别为 $0, \pm\mu_{\text{B}}B/h$。
(c) 计算分裂量，分别以 GHz 和波长的皮米数表示。
(d) 多数实际谱线分裂为多于三个的
分量（“反常”\omterm{prop:b3:quantum-angular-momentum:magnetic}{塞曼效应}）：由电子自身携带的哪一种缺失成分
曾为其提供历史证据？
\end{exercise}

\begin{exercise}[$\star\star$]\label{exo:b3:quantum-angular-momentum:8}
\omterm{thm:b3:quantum-angular-momentum:radial}{离心势垒}。对氢原子（$V = -e^2/4\pi\varepsilon_0r$），
写出 $U_{\text{有效}}$，取 $l = 1$，并以 $E_{\text{I}}$ 和
$a_0$ 为单位。(a) 求其极小值点及极小值。(b) 证明 $U_{\text{有效}} >
0$，当 $r < \dots$ 时成立——求出势垒占主导的
半径。(c) 比较 $|\psi|^2$ 在 $r = 0$ 附近，对 $s$ 态与
$p$ 态（$u \sim r^{l+1}$，在本阶段不加证明地承认）：哪些态“触及”
原子核？(d) 电子俘获——原子核吞并其原子中的一个
电子——几乎完全发生于 $s$ 壳层：
用一句话解释。
\end{exercise}

\begin{exercise}[$\star\star$]\label{exo:b3:quantum-angular-momentum:9}
在 $\ket{l, m}$ 上：(a) 证明 $\langle\hat L_x\rangle = \langle\hat
L_y\rangle = 0$；(b) 证明 $\langle\hat L_x^2\rangle = \langle\hat
L_y^2\rangle = \tfrac12[l(l+1) - m^2]\hbar^2$；(c) 验证
\omterm{thm:b3:quantum-formalism:uncertainty}{不确定性关系} $\Delta L_x\,\Delta L_y \ge \tfrac\hbar2
|\langle\hat L_z\rangle|$，在最大投影态 $m = l$ 上；(d) 解释
图中“锥面”图像，并结合
(a) 与 (b)。
\end{exercise}

\begin{exercise}[$\star\star\star$]\label{exo:b3:quantum-angular-momentum:10}
振动–转动带。一次红外振动跃迁
（$\hbar\omega_0$）同时使双原子分子的 $J$ 改变 $\pm1$。
(a) 证明吸收线位于 $h\nu =
\hbar\omega_0 + 2B(J+1)$（$R$ 支，由 $J \to J+1$）和
$h\nu = \hbar\omega_0 - 2BJ$（$P$ 支，由 $J \to J-1$），$J \ge
1$。(b) 为何\emph{空隙}恰在 $\hbar\omega_0$ 处出现？(c)
画出该带：两列梳状谱线夹着缺失的中央线，
每条谱线的强度由
\cref{exo:b3:quantum-angular-momentum:6}给出。(d)
\cref{pb:b3:harmonic-oscillator:1}中的 \qty{15}{\micro m} 带恰具有这一结构：
究竟什么决定其总宽度，进而
决定使新增二氧化碳有效的“翼”？
\end{exercise}

\begin{exercise}[$\star\star\star$]\label{exo:b3:quantum-angular-momentum:11}
完成代数推导。(a) 由 $\hat J_\mp\hat J_\pm = \hat J^2 -
\hat J_z^2 \mp \hbar\hat J_z$，计算 $\|\hat J_\pm\ket{j,m}\|^2$
及升降系数。(b) 证明：若升降链未能终止，
则某一范数将为负：给出该态。
(c) 由此得出 $m_{\max} - m_{\min} = 2j$ 必为非负
整数，并列出允许的 $j$。(d) 论证究竟在何处
\emph{失效}，未能排除半整数——而哪一
附加要求（球面上的单值性）排除了它们，
且仅对轨道运动成立？
\end{exercise}

\begin{exercise}[$\star\star\star$]\label{exo:b3:quantum-angular-momentum:12}
爱因斯坦–德哈斯：角动量是力学量。一个铁制圆柱体
（半径 $a = \qty{1.0}{cm}$，密度 $\rho = \qty{7.9e3}{kg/m^3}$，
$n = \qty{8.5e28}{\text{原子}/m^3}$）悬挂于一根扭丝上，翻转
其磁化方向时，每个原子约翻转一个 $\hbar$ 的
角动量。(a) 由守恒，杆必定反冲：计算
单位体积所转移的角动量。(b) 已知杆的转动惯量 $\tfrac12 M
a^2$，证明杆获得 $\omega = 2n\hbar/\rho a^2$，
并求其数值。(c) 1915 年的实验在扭丝的
共振频率下驱动翻转，以累积微小冲量：估计
经 $Q \sim 100$ 次共振翻转后的振动振幅，并取每次冲量对应的
$\omega$。(d) 实测的磁矩与角动量之比，
恰为轨道预言值 $e/2m_{\text{e}}$ 的两倍：这
个因子 2 预示着什么？
\end{exercise}

\section{问题：在 2.6 毫米处称量星系}

\begin{problem}[{周末问题——一氧化碳、射电
天文与寒冷宇宙}]\label{pb:b3:quantum-angular-momentum:1}
星系中大部分成星物质是冷的分子氢——
而它不可见：H$_2$ 呈对称结构，无偶极矩，且
在冷云温度下无转动谱线。天文学的解法在于
其忠实的伴侣。每 $\num{e4}$ 个 H$_2$ 分子对应一个 CO 分子，
其 \qty{115}{GHz} 的转动谱线阶梯，点亮
天空中每一块分子云。数据：对 CO，$B/h = \qty{57.6}{GHz}$（$B =
\qty{0.238}{meV}$）；键长 $r_0 = \qty{0.113}{nm}$；约化
质量 $\mu = \qty{1.14e-26}{kg}$；$k_{\text{B}}T$ 在 \qty{10}{K} 时为
\qty{0.862}{meV}。

\medskip\noindent\textbf{第一部分——量子转子。}
\begin{enumerate}
  \item 由 $\hat H = \hat L^2/2I$ 给出能量 $E_J$ 及
        其简并度。
  \item 验证 CO 的 $B = \hbar^2/2I$（由 $r_0$ 与 $\mu$ 出发）。
  \item 利用选择定则 $\Delta J = \pm1$，导出
        吸收频率 $\nu_{J+1\leftarrow J} = 2B(J+1)/h$，
        并列出前三条谱线。
  \item 谱线为何\emph{等间距}？所测间距
        通过 $I$ 能向天文学家提供什么？
  \item 计算基频谱线的波长，并解释为何观测它需要
        高而干燥的观测站址（什么吸收
        地球大气中的毫米波——回想
        \cref{pb:b3:harmonic-oscillator:1}中那个分子的
        近邻，水）？
  \item $1 \to 0$ 谱线的光子携带
        $\hbar$ 的角动量：为何 $\Delta J = \pm1$ 不仅仅是
        一条经验法则，而是一条守恒定律？
\end{enumerate}

\medskip\noindent\textbf{第二部分——为何用 CO 而不用 H$_2$。}
\begin{enumerate}[resume]
  \item H$_2$ 是同核分子：说明它为何不发出任何转动
        偶极辐射。
  \item H$_2$ 还很\emph{轻}：计算 H$_2$ 的转动常数
        （$\mu = m_{\text{p}}/2$，$r_0 = \qty{0.074}{nm}$）
        及其最低可及跃迁能量，并与
        \qty{10}{K} 时的 $k_{\text{B}}T$ 比较：即便借助四极
        “后门”跃迁，冷的 H$_2$ 能否辐射？
  \item CO 的第一级阶梯需 $2B = \qty{0.48}{meV}$，对比
        $k_{\text{B}}T = \qty{0.86}{meV}$：在
        \qty{10}{K} 时，该阶梯是否仍然有效？
  \item 计算 \qty{10}{K} 时的布居数之比 $n_1/n_0$
        （已计入简并度）。
  \item \qty{10}{K} 时布居数在哪个 $J$ 附近达到峰值
        （取 $J_{\max} \approx \sqrt{k_{\text{B}}T/2B} - \tfrac12$）？
  \item 用一句话概括这笔“交易”的实质：CO 提供了什么，
        以及必须假设 CO 与 H$_2$ 之间的
        比值是多少。
\end{enumerate}

\medskip\noindent\textbf{第三部分——读取天空。}
\begin{enumerate}[resume]
  \item 某分子云的谱线在 \qty{115.156}{GHz} 处到达，而非
        实验室中的 \qty{115.271}{GHz}：计算该分子云沿视线的速度
        及其方向。
  \item 谱线因内部运动产生多普勒\emph{展宽}，
        内部运动的速度为 $\pm\qty{2}{km/s}$：以 MHz 计算线宽，并
        与 \qty{10}{K} 时预期的热展宽相比较
        （$v_{\text{th}} \sim \sqrt{k_{\text{B}}T/m_{\text{CO}}}
        \approx \qty{55}{m/s}$）：何者主导展宽？
  \item 沿旋涡星系的盘面映射 CO 的多普勒频移，
        得到其旋转曲线 $v(r)$。这些曲线在远超可见盘面的
        范围内仍保持\emph{平坦}。回想（本科一年级分册，
        万有引力）中围绕中心集中质量的 $v(r)$ 应如何变化，并
        说明这种平坦性
        意味着什么。
  \item $2\to1$ 谱线与
        $1\to0$ 谱线的亮度之比，用于测量能级布居：它给出
        该分子云的哪一个单一物理量？
  \item ALMA 能分辨光在宇宙年轻时离开的星系中的 CO；
        所观测的
        $1 \to 0$ 谱线频率，星系位于“红移 $z = 1$”处，
        恰为静止值的一半：该频率为多少，并
        （回想\cref{ex:b3:relativistic-kinematics:redshift}）
        是什么拉伸了它？
  \item 由测得的 CO 谱线光度，天文学家给出分子云的
        质量（以太阳质量计）：列出换算链（谱线
        光子 $\to$ CO 数 $\to$ H$_2$ 质量），以及其中最薄弱的
        环节。
\end{enumerate}

\medskip\noindent\textbf{第四部分——寒冷的摇篮。}
\begin{enumerate}[resume]
  \item 恒星形成核心处于 \qty{10}{K}：由维恩位移定律，
        其热辐射峰值位于 \qty{290}{\micro m} 附近——
        所有光学望远镜都看不见。用一句话说明：为什么
        转动能级梯子是这样一个核心所能提供的\emph{唯一}
        温度计与标尺？
  \item $J = 1$ 能级位于基态之上 \qty{5.5}{K}（以温度单位计）：
        为什么这使得 $1\to0$ 谱线
        恰好在 \qty{10}{K}
        区间内成为极佳的温度计（而不是，比如说，位于
        \qty{2}{eV} $\sim \qty{23000}{K}$ 的光学谱线）？
  \item 为什么\emph{坍缩中的}核心最终会在 CO（$1\to0$）中
        不再可见——考虑高密度与
        尘埃对谱线及其逃逸的作用。
  \item 分子云还显示 NH$_3$、HCN 以及数十种
        其他分子的谱线，每一种都有自己的 $B$ 与偶极矩：
        这种“化学无线电调频盘”的丰富性能让天文学家
        分辨出什么？
  \item 整个天空以 \qty{2.7}{K} 发光（宇宙微
        波背景）：云团 CO 谱线的对比度，
        随云团自身温度趋近 \qty{2.7}{K}，将发生什么变化？
        以及为什么这对该温度计而言是一道
        硬性下限？
  \item 毫米波天线的反射面必须保持其形状至约
        $\lambda/20$：在 \qty{115}{GHz} 下计算该容差，
        并与光学反射镜（\qty{500}{nm}）必须满足的容差相比较——
        为什么 ALMA 的十二米天线
        能够在露天环境中建造？
  \item 概括这一著名结论：能级梯子 $BJ(J+1)$，其中 $B/h =
        \qty{57.6}{GHz}$，在十开尔文下布居，并在
        \qty{2.6}{mm} 处读取，正是测定冷宇宙的质量、温度、
        速度与转动——以及螺旋星系周围暗物质的证据——
        所采用的方法。
\end{enumerate}
\end{problem}
